\documentclass[11pt]{article}
\usepackage[T1]{fontenc}
\usepackage[textwidth=250pt]{geometry}
\usepackage{amsmath,amssymb}
\usepackage{microtype}
\pagestyle{empty}
\begin{document}
Let $f(x)=\sin x+\cos x$ on the interval $[0,2\pi]$, and let $g(x)=x^3-1$.
Then $f(0)=1$ and $g(1)=0$, so the composition $f(g(1))=1$ as well; the
limit $\lim_{x\to 0}\frac{\sin x}{x}=1$ is used with $\log 2<1$ in part (c).

The map $h\colon A\to B$ with $h(a)=a+1$ is a bijection when $A=B=\mathbb{Z}$,
and its inverse is $h^{-1}(b)=b-1$; the set $\{1,2,3\}$ is sent to the set
$\{2,3,4\}$, and the function $\max(a,b)$ is the larger of the two values.

The spacing commands give $a\,b$ and $a\quad b$ and $x\;y$ in one line, and
the text keeps going after them with the formula $p\mid q$ and then $n!=120$
for $n=5$, which finishes the list of examples in this last paragraph.
\end{document}
